\documentclass[10pt, a4paper, twocolumn]{article}
\usepackage{kotex}
\usepackage[margin=18mm]{geometry}
\usepackage{amsmath, amssymb, amsfonts, bm}
\usepackage{booktabs}
\usepackage{hyperref}
\usepackage{xcolor}
\usepackage{microtype}

\hypersetup{
    colorlinks=true,
    linkcolor=blue!70!black,
    citecolor=blue!70!black,
    urlcolor=blue!70!black
}

\title{\LARGE \textbf{[연구제안서] Multimodal Long Video Token Compression}\\[0.3em]
\large 시공간 2축 적응형 토큰 압축 기반 차세대 온디바이스 옴니모달 아키텍처}
\author{\textbf{정현우 (Hyeonwoo Jung)}\\[0.2em]
POSTECH / AI 대학원 지원 연구계획서\\
\texttt{hyeonwoo.research@postech.ac.kr}}
\date{2026년 08월 26일}

\begin{document}
\maketitle

\begin{abstract}
최신 멀티모달 대형 언어 모델(MLLM)은 고해상도 영상 및 오디오 처리 시 트랜스포머 Self-Attention의 이차 복잡도 $O(N^2)$로 인해 극심한 연산 및 메모리 절벽에 직면한다. 90분 분량의 비디오는 최대 5,400만($54\text{M}$) 토큰을 유발하여 온디바이스 실시간 추론을 마비시킨다. 본 연구는 선행 연구인 FastV(깊이 축 토큰 프루닝), StreamKV(시간 축 캐시 관리), OmniSelect(질문 기반 모달리티 게이팅)의 한계를 심층 분석하고, 이들의 모순을 해소하는 \textbf{시공간 2축 적응형 압축 아키텍처(DA-MLLM)}를 제안한다. 제안 기법은 FLOPs 65\% 절감, KV Cache 메모리 80\% 감축, 초당 프레임 3.2배 향상을 달성하면서도 Video-MME 등 핵심 벤치마크 성능을 98.5\% 이상 방어하는 것을 목표로 한다.
\end{abstract}

\section{연구 배경 및 극한 병목 (Problem)}
트랜스포머의 Self-Attention 메커니즘은 시퀀스 길이 $N$에 대해 $O(N^2)$ 연산 복잡도와 $O(N)$ KV Cache 메모리 풋프린트를 요구한다.
\begin{itemize}
    \item \textbf{토큰 팽창 수치}: 텍스트 프롬프트가 수백 토큰인 반면, 90분 영상은 약 5,400만($54\text{M}$) 토큰을 생성.
    \item \textbf{온디바이스 한계}: 스마트 글래스 및 로봇 엣지 칩셋(Jetson, NPU)의 한정된 SRAM 및 대역폭에서 실시간 서빙 불가능.
\end{itemize}

\section{선행 연구 맹점 비판 (Prior Art Critique)}
기존 연구들은 특정 모달리티나 단일 축에 치우쳐 상호 충돌을 겪는다:
\begin{itemize}
    \item \textbf{FastV (ECCV 2024)}: Layer 2 이후 하위 50\% 토큰 드롭으로 FLOPs 45\% 절감했으나, 정적 이미지에 국한되어 비디오 시간축 인과관계(Temporal Causality) 왜곡 및 정적 고정 비율 컷으로 미세 정보 소실 발생.
    \item \textbf{StreamKV}: 슬라이딩 윈도우로 KV Cache는 줄이나 토큰 텐서 자체의 FFN 연산은 감축 불가.
    \item \textbf{OmniSelect (2026)}: 질문에 따라 모달리티 프루닝 비율을 동적 조절하나, 트랜스포머 내부 레이어 깊이별 정보 집약(Anchor Sink) 특성을 반영하지 못함.
\end{itemize}

\section{핵심 제안 기법 (Proposed Method)}
본 제안서는 \textbf{Dual-Axis Adaptive MLLM (DA-MLLM)}을 통해 깊이 축과 시간 축을 동시 압축한다.

\subsection{수학적 수식화}
1. \textbf{질문 반응형 모달리티 가중치 산출}:
\begin{equation}
\alpha_v, \alpha_a = \text{Softmax}(\mathbf{W}_g \cdot [\text{Embed}(Q); \text{Pooling}(V); \text{Pooling}(A)])
\end{equation}
2. \textbf{레이어 축 깊이별 적응형 프루닝}:
질문 복잡도 $\mathcal{C}(Q)$에 따라 프루닝 시작 레이어 $K^*$와 제거율 $R^*$를 동적 결정:
\begin{equation}
K^* = \lfloor K_0 + \beta \cdot \mathcal{C}(Q) \rfloor, \quad R^* = R_0 \cdot (1 - \alpha_v)
\end{equation}
Layer $K^*$ 통과 후 평균 어텐션 점수 $\phi_{\text{attn}}(t)$ 하위 $R^*$ 토큰을 FFN 통과 직전 텐서에서 영구 제거.

3. \textbf{시간 축 캐시 슬라이딩 정렬}:
가변 토큰 프루닝 시 발생하는 PagedAttention 메모리 인덱싱 충돌을 해결하기 위해, \textbf{Time-Anchor Indexing Mask}를 설계하여 공간 토큰은 삭제하되 각 타임스탬프의 기준 앵커 토큰만 KV 캐시에 압축 유지.

\section{정량적 목표 및 실험 검증}
\begin{itemize}
    \item \textbf{연산량(FLOPs)}: 베이스라인 대비 \textbf{65\% 이상 절감}.
    \item \textbf{메모리 풋프린트}: 1시간 비디오 스트리밍 시 KV Cache \textbf{80\% 이상 감축}.
    \item \textbf{추론 속도}: Jetson Orin 엣지 기준 \textbf{3.2배 향상} (15 FPS $\rightarrow$ 48 FPS 실시간성 확보).
    \item \textbf{정확도 보존}: Video-MME, MMMU 벤치마크 \textbf{98.5\% 이상 유지}.
\end{itemize}

\section{킬러 활용처: 온디바이스 라이프로깅}
사용자의 시야와 음성을 24시간 실시간 캡처하는 스마트 글래스(AI Life-Companion) 환경에서, 배터리 방전과 발열 없이 초장시간 비디오 인과 추론 완벽 구현.

\end{document}
